\subsection{Removing source core three-handles relative to the bundle ports}
\label{marked-cut:three-handles}

The source core retains one three-handle for each interior original vertex.
These are vertex-link balls, not compressed families. Its two-handles are
the extreme gaps of interior source edges. This gives a finite dual graph
in which the boundary is a distinguished root. We can contract a forest
from the interior vertices to this root before testing one/two cancellation.

\begin{proposition}[Authenticated dual endpoints]
\label{marked-cut:dual-endpoints}
Let an interior source edge $e$ have normal weight $w$, with its canonical
local endpoints ordered as $v_0,v_1$. If $w=0$, the core has one two-handle,
whose two dual endpoints are the vertex-link regions of $v_0,v_1$. If
$w>0$, the rank-zero gap has endpoints $v_0$ and the cut boundary, while
the rank-$w$ gap has endpoints on the cut boundary and $v_1$.
A source endpoint on the original boundary is also a boundary endpoint.
Repeated interior-vertex endpoints count twice.
\end{proposition}
\begin{proof}
The dual cocore of an edge-gap two-handle is the source gap interval itself,
transverse to the edge-link disk. With zero weight that interval extends
between the two original vertices. With positive weight, its two extreme
gaps extend from the corresponding vertex to the first or last normal-disk
intersection, respectively. The latter endpoints are exposed by the cut.
A vertex contributes a three-handle precisely when its link is a sphere;
a boundary-vertex link is already a boundary disk. Thus the stated endpoint
list is the geometric incidence list with three-handle attaching spheres.
For a loop edge with both ends at the same interior vertex, the two distinct
ends of the cocore give two intersections, regardless of the equal label.
\end{proof}

\begin{proposition}[Complete dual boundary forest]
\label{marked-cut:dual-forest}
For the validated connected compact source with nonempty boundary, all core
three-handles can be removed by $O(T)$ two/three-handle cancellations. The
source zero/one handles and all bundle port charts are preserved. Only the
cancelled two-handles and three-handles are deleted.
\end{proposition}
\begin{proof}
Create one dual vertex for each interior source vertex and identify all
original and cut-boundary endpoints with a root. The original source
one-skeleton is connected and contains a boundary vertex. A path from an
interior vertex towards the original boundary either follows zero-weight
edges until reaching that boundary, or first encounters a positive-weight
edge. In the second case its corresponding extreme gap already gives an
edge to the cut boundary. Every dual vertex therefore reaches the root.
Source edges touching an interior vertex have circular links, so their
required extreme-gap two-handles are all present in the core.

Starting at the root, choose an edge to a remaining interior vertex. Its
two-handle belt sphere meets exactly that remaining three-handle attaching
sphere once. The geometric two/three cancellation removes this pair. All
other endpoints at the removed vertex become boundary endpoints. This is
the usual dual boundary-forest contraction, and the argument repeats until
no three-handle remains. Loops at a remaining vertex have two intersections
and cannot supply the pivot. There are at most $4T$ interior vertices.

The two/three cancellation changes the top handles, leaving all other
attaching maps into the original zero/one body unchanged. The bundle ports
lie in that lower body, outside the top-handle attaching regions. The
cancellation can therefore be performed relative to the marked ports.
Deleting equal numbers of two- and three-handles also preserves the core
Euler characteristic used in the independent whole-cut consistency check.
\end{proof}

\paragraph{Independent source replay.}
The producer derives endpoint owners from the source vertex partition and
normal edge weights. The independent checker instead assigns local vertex
occurrences directly from the authenticated three-handle link triangles.
It computes each root-edge weight directly from the local normal coordinates:
triangles at its two ends plus quadrilaterals that meet that edge. No
producer geometry or cancellation planner is needed. A record supplies the
complete dual endpoint list, every cancelled pair, and the ordered list of
surviving original two-handle identifiers. The checker preserves endpoint
multiplicity, checks a single remaining endpoint at each move, and renumbers
only the surviving two-handles. The boundary kernel then runs on that
verified reduced core with the original profiles and port words.

\paragraph{Sparse algorithm and accounting.}
The producer indexes both occurrences of each dual edge endpoint. A queue
holds edges with one live interior endpoint; its degree and xor of live
endpoint labels identify the cancellable pair. Removing a vertex scans each
of its incident endpoint occurrences once. A self-loop is initially degree
two and is never an eligible pivot; a temporarily queued record created
while deleting its two occurrences is stale and is skipped. Each dual edge
has at most two endpoint occurrences, so this phase has $O(T)$ indexed work
and storage. The checker simply filters the two authenticated endpoints
against its live-vertex set at each recorded move. Full coverage and the
ordered surviving-row list are checked independently. Source preparation,
integer reading, and the existing whole-cut orbit queries have their own
costs; this is not a linear bound on the complete cut operator.

The input contract is the source-authenticated parallelity geometry. The
producer uses its ordinary vertex partition and edge-weight analysis. The
checker imports neither of those producer routines, neither the source
geometry builder nor the cancellation queue. It reconstructs owners from
the authenticated link triangles and computes the root-edge weight locally.
Equal numbers of deleted two- and three-handles preserve the original
summary's Euler characteristic, so the unchanged independent genus identity
still applies.

\paragraph{Whole-cut integration and ordering.}
New proofs use outer schema \path{normal-cut-handlebodies-v6}. They contain
the original parallelity and boundary proofs, a
\path{core-three-handle-reduction-v1} record, and the existing timestamp-based
boundary-kernel proof. The producer and checker pass the reduced geometry to
that kernel; source proofs continue to authenticate the original geometry.
Only two-handle identifiers are renumbered, in surviving source order.
Core zero/one handles and every port chart retain their original identifiers.
Sources without three-handles keep their previous certificate routes.

The reduction must precede the source-core stopping-set precheck. A stopping
subset of original rows can disappear when two/three cancellation deletes
some of those rows. Applying the precheck to the unreduced core would thus
be unsound as a completeness-preserving filter for this enlarged method.
After reduction, the previous monotonicity proof applies to the remaining
core rows. Exhausting an orbit allowance or interrupting a callback still
emits no partial certificate.

This extends the existing source-handle cancellation argument of
Section~\ref{sec:source-handles} to arbitrary supplied normal cuts by using
extreme edge gaps and cut-boundary endpoints. The local cancellation rule is
Procedure 6 in Section 6.1 of~\cite{lackenby-parallelity}; the source argument
above establishes its specific charts and preservation of the bundle ports.

\paragraph{Source tests and present scope.}
All 33 focused tests pass in 73.185 seconds. They cover all 106 historical
source cores, fourteen additional triangulated sources, nested interior
vertices, a genuine interior loop, malformed records, source binding,
independent replay with producers disabled, binary multiplicity $2^{512}$,
relabelled sources, callback interruption, and whole-query budgets.

The additional fixtures are regenerated by topology-preserving Regina
Pachner moves. Repeated $1$--$4$ subdivisions select the tetrahedron with the
fewest boundary vertices, producing up to twelve interior vertices in 37
tetrahedra. The deeper examples require cancellations along interior dual
edges, after their parents reach the boundary. A $2$--$3$ move after one
subdivision creates an edge whose two ends are the same interior vertex.
The checker rejects this edge as an initial two/three pivot.

All fourteen sources admit complete three-handle reduction. Ten also yield
whole-cut certificates, and all ten independently agree with explicit
Regina cuts. Four remain inconclusive. Two cut off interior vertex spheres
and leave non-handlebody components according to Regina. The other two,
the empty and boundary-link cuts of the loop fixture, are genuine
handlebody censuses missed by this chosen raw core presentation. Thus
complete three-handle elimination is not complete handlebody recognition.
The fixture generator, exported pairings and coordinates, independent
oracle records, and source pins are retained.

The fresh historical audit certifies 92 of 106 supplied cuts, up from 85.
All 92 independently replay, 85 agree with literal Regina cuts, and seven
large-coordinate cases exceed the expansion cap. The fourteen remaining
refusals comprise eight raw core stalls, three square core obstructions,
and three joint boundary-kernel stalls. None is a topological negative
certificate. The seven newly certified cuts are the three interior layered
torus pilots, the interior trefoil pilot, the mixed interior $(1,2)$ and
$(3,5)$ tori, and the mixed interior figure-eight exterior.

The historical compatibility comparison, full regression suite and paired
performance measurements are recorded below when complete.


\paragraph{Dependence on the selected dual forest.}
For each of the two loop misses, a separate source-bound experiment enumerates
all four eligible first pairs. The native queue selects the one that leaves
a blocked raw matrix. Each of the other three choices yields a complete v6
certificate, and all six alternative certificates independently replay.
The records in \path{data/three-handles-choice-probe.json} are an experiment,
not the production selector. They show that forest choice, rather than a
lack of valid certificates in this format, explains these two misses.
Choosing a useful forest in general needs a separate selection algorithm
and complexity argument; merely enumerating all spanning forests can be
exponential. The current queue's completeness concerns elimination of
three-handles, irrespective of the later raw cancellation outcome.

